\documentclass[11pt, twocolumn, landscape, a4paper]{article}

\usepackage[margin=1cm]{geometry}
\usepackage[utf8]{inputenc}
\usepackage[spanish]{babel}
\usepackage{amssymb, amsmath, amsbsy}
\usepackage{tasks}
\usepackage{enumerate}
\usepackage[pdftex]{hyperref}
\hypersetup{pdftitle={Aritmética - Varios},pdfauthor={Luis Carrillo Gutiérrez}}

\fontfamily{cmss}\selectfont
\renewcommand{\familydefault}{\sfdefault}
\pagestyle{empty}

\begin{document}
\subsection*{Aritmética - conjuntos}
\begin{itemize}
\item{Hallar el cardinal del conjunto siguiente para indicar el número de subconjuntos ternarios que tiene : $B = \{x \in \mathbb{Z}^{+}\,/\,(x > 8) \to (x = 2)\}$ 
\begin{tasks}(5)
\task{42}\task{45}\task{48}\task{56}\task{63}
\end{tasks}
}
\item{Si $n[P(A)] = 128$; $n[P(B)] = 32$; $n[P(A \cap B)] = 8;$ : Hallar el cardinal de $P(A \cup B)$ sumado con el cardinal de : $C = \Big\{(3x + 1) \in \mathbb{Z}^{+}\,/\,x < \dfrac{5}{3}\Big\}$
\begin{tasks}(5)
\task{512}\task{517}\task{519}\task{520}\task{521}
\end{tasks}
}
\item{Oscar compra $9$ baldes de pinturas de diferentes colores. Los mezcla en igual proporción. ¿Cuántos nuevos matices se pueden obtener?
\begin{tasks}(5)
\task{246}\task{247}\task{502}\task{503}\task{512}
\end{tasks}
}
\item{El conjunto A tiene 200 subconjuntos no ternarios. ¿Cuántos subconjuntos quinarios tendrá?
\begin{tasks}(5)
\task{48}\task{56}\task{64}\task{72}\task{81}
\end{tasks}
}
\end{itemize}
\end{document}